\documentclass[11pt]{article}
\usepackage[a4paper,margin=28mm]{geometry}
\usepackage{amsmath,amssymb,amsthm}
\usepackage[colorlinks=true,linkcolor=blue,urlcolor=blue,citecolor=blue]{hyperref}
\newtheorem{theorem}{Theorem}
\newtheorem{lemma}{Lemma}
\newtheorem{corollary}[lemma]{Corollary}
\newtheorem{proposition}[lemma]{Proposition}
\newcommand{\Q}{\mathbb Q}
\newcommand{\Z}{\mathbb Z}
\newcommand{\F}{\mathbb F}
\DeclareMathOperator{\End}{End}
\DeclareMathOperator{\Gal}{Gal}
\DeclareMathOperator{\GL}{GL}
\DeclareMathOperator{\PGL}{PGL}
\DeclareMathOperator{\Hom}{Hom}
\DeclareMathOperator{\Aut}{Aut}
\DeclareMathOperator{\Ad}{Ad}
\title{An obstruction to rational \(63\)-torsion on geometrically split abelian surfaces}
\author{}
\date{21 October 2026}

\begin{document}
\maketitle

\begin{theorem}\label{thm:main}
Let \(A/\Q\) be an abelian surface that is simple over \(\Q\) and
geometrically split. Then \(A(\Q)\) has no point of order \(63\).
\end{theorem}

Here ``geometrically split'' means that \(A_{\overline{\Q}}\) is
isogenous to a product of elliptic curves. The theorem places no
restriction on the degree of this isogeny, on the field of definition of
the geometric endomorphisms, on the integral endomorphism ring, or on the
polarization; in particular, it does not require a principal
polarization. Consequently, the group
\href{https://genus-2-torsion.github.io/group.html?g=63}{\([63]\)}
cannot occur as the rational torsion subgroup of a geometrically split
genus-\(2\) Jacobian that is simple over \(\Q\).

The hypothesis of simplicity over \(\Q\) is essential. A product of
elliptic curves over \(\Q\) with rational points of orders \(7\) and
\(9\) has a rational point of order \(63\), and the census contains a
genus-\(2\) Jacobian that is split over \(\Q\) and has torsion subgroup
\([63]\).

\section{Inputs}\label{sec:inputs}

We use the following results on torsion points of elliptic curves.
\begin{enumerate}
\item Kenku--Momose \cite[Theorem (0.1)]{KM}: let \(N\) be a product of
powers of \(2,3,5,7,11,13\), and let \(m\mid N\). If \(X_1(m,N)\) has
positive genus and
\[
 (m,N)\notin\{(1,11),(1,13),(1,14),(1,15),(1,16),(1,18),(2,10),(2,12)\},
\]
then \(Y_1(m,N)(k)=\emptyset\) for every quadratic field \(k\). The
curve \(X_1(21)=X_1(1,21)\) has genus \(5\). Hence \emph{no elliptic
curve over a quadratic field has a point of order \(21\)}.
\item Mazur \cite[Theorem (7\('\))]{Mazur}: if an elliptic curve over
\(\Q\) has a rational point of order \(m\), then \(m\leq10\) or
\(m=12\). In particular \(Y_1(21)(\Q)=\emptyset\).
\item Derickx--Etropolski--van Hoeij--Morrow--Zureick-Brown
\cite[Theorem A]{DEvHMZB} classify the torsion of elliptic curves over
cubic fields. Theorem A states in addition that the elliptic curve with
Cremona label 162b1,
\[
 y^2+xy+y=x^3-x^2-5x+5,
\]
over \(\Q(\zeta_9)^+=\Q[x]/(x^3-3x+1)\) is the unique elliptic curve
over a cubic field with torsion \(\Z/21\Z\). We use the precise
statement proved in \cite[Section 6.2]{DEvHMZB}: the only closed points
of degree \(3\) on \(X_1(21)\) are two points \(D_0\) and \(D_0'\).
\end{enumerate}

For the reduction of abelian varieties we use the following results of
Serre--Tate \cite{ST}. Let \(K\) be a finite extension of \(\Q_p\) with
residue field \(k\), let \(A\) be an abelian variety over \(K\), and let
\(\ell\neq p\) be a prime.
\begin{enumerate}
\item[(a)] (N\'eron--Ogg--Shafarevich, \cite[Theorem 1]{ST}.) \(A\) has
good reduction if and only if the inertia group of \(K\) acts trivially
on \(T_\ell(A)\). By \cite[Corollary 2 of Theorem 1]{ST}, two
\(K\)-isogenous abelian varieties both have, or both do not have, good
reduction.
\item[(b)] (\cite[Lemma 2]{ST}.) For \(m\) prime to \(p\), reduction is
an isomorphism from the inertia invariants of \(A[m]\) onto the
\(m\)-torsion of the special fibre of the N\'eron model, compatible with
the Galois action. Hence, if \(A\) has good reduction \(\widetilde A\),
then \(A(K)[m]\) injects into \(\widetilde A(k)\), and
\(T_\ell(A)\cong T_\ell(\widetilde A)\) compatibly with Frobenius.
\item[(c)] (\cite[\S2, p.~498]{ST}.) If \(A\) has potentially good
reduction, then there is a finite extension \(K'/K\) with residue field
\(k\) over which \(A\) has good reduction.
\end{enumerate}

For an abelian surface \(B\) over \(\F_q\), \(\#B(\F_q)=P(1)\), where
\(P\) is the characteristic polynomial of Frobenius, and the complex
roots of \(P\) have absolute value \(\sqrt q\) (Weil). Hence
\(\#B(\F_q)\leq(1+\sqrt q)^4\).

We also use the following standard facts: the independence of the
semisimplified residual representation from the choice of a stable
lattice (Brauer--Nesbitt), which applies to isomorphic \(p\)-adic
representations after transporting a lattice through the isomorphism;
the Skolem--Noether theorem; Poincar\'e's complete reducibility theorem;
Galois descent for twists of projective varieties by a cocycle with
values in the automorphism group; the Weil pairing; and the existence of
a polarization defined over \(\Q\) on every abelian variety over \(\Q\).

\emph{Notation.} For a number field \(K\subset\overline{\Q}\), put
\(G_K=\Gal(\overline{\Q}/K)\). For a prime \(\ell\), \(T_\ell\) and
\(V_\ell=T_\ell\otimes\Q_\ell\) denote the Tate module and the rational
Tate module. For a stable lattice \(T'\) in a \(p\)-adic representation,
\(T'/pT'\) is its residual module and \((\cdot)^{\mathrm{ss}}\) denotes
semisimplification. We write \(\overline\chi_p\) for the mod-\(p\)
cyclotomic character. A rational point of order \(p\) spans a trivial
submodule of \(A[p]=T_p(A)/pT_p(A)\); by lattice independence, the
residual module of every \(G_{\Q}\)-stable lattice in \(V_p(A)\) then has
a trivial composition factor.

\section{Setup and two lemmas on elliptic curves}\label{sec:setup}

Throughout the proof, \(A/\Q\) is an abelian surface that is simple over
\(\Q\) and geometrically split, and we suppose that \(P\in A(\Q)\) has
order \(63\). Put
\[
 P_7=9P,\qquad P_3=21P,
\]
of orders \(7\) and \(3\). Exactly one of the following holds:
\begin{itemize}
\item[(N)] \(A_{\overline{\Q}}\) is isogenous to \(E_1\times E_2\),
where \(E_1\) and \(E_2\) are not isogenous;
\item[(CM)] \(A_{\overline{\Q}}\) is isogenous to \(E^2\), where \(E\)
has complex multiplication;
\item[(S)] \(A_{\overline{\Q}}\) is isogenous to \(E^2\), where \(E\)
has no complex multiplication.
\end{itemize}

\begin{lemma}[Combining coprime torsion]\label{lem:combine}
Let \(m,n\) be coprime positive integers and \(K\) a number field.
Suppose \(E_m\) and \(E_n\) are \(K\)-isogenous elliptic curves over
\(K\) with points of orders \(m\) and \(n\) in \(E_m(K)\) and
\(E_n(K)\), respectively. Then their \(K\)-isogeny class contains an
elliptic curve with a \(K\)-rational point of order \(mn\).
\end{lemma}

\begin{proof}
Let \(f:E_m\to E_n\) be a \(K\)-isogeny, and let \(H\) be the subgroup
of \(\ker f\) consisting of the elements whose order is divisible only
by primes dividing \(n\). It is a characteristic subgroup of \(\ker f\),
hence \(G_K\)-stable, and \(f\) factors over \(K\) as
\[
 E_m\xrightarrow{\ a\ }E'=E_m/H\xrightarrow{\ b\ }E_n .
\]
The degree of \(a\) is prime to \(m\), so \(a\) is injective on
\(E_m[m]\) and maps the given point of order \(m\) to a point of order
\(m\) in \(E'(K)\). The degree of \(b\) is prime to \(n\); since
\(b\circ\hat b=[\deg b]\) is injective on \(E_n[n]\), so is the dual
isogeny \(\hat b\), which maps the given point of order \(n\) to a point
of order \(n\) in \(E'(K)\). The sum of the two points has order
\(mn\).
\end{proof}

\begin{lemma}\label{lem:sequence}
Let \(X\) be an abelian surface over a field \(k\), let \(E\subset X\)
be an elliptic curve defined over \(k\), and put \(F=X/E\). If
\(Q\in X(k)\) has prime order \(q\), then either \(Q\in E(k)\), or the
image of \(Q\) in \(F(k)\) has order \(q\).
\end{lemma}

\begin{proof}
The image has order dividing \(q\). If it is zero, then
\(Q\in E(\overline k)\cap X(k)=E(k)\).
\end{proof}

\section{Nonisogenous factors}\label{sec:nonisogenous}

Suppose that \(A_{\overline{\Q}}\) is isogenous to \(E_1\times E_2\),
where \(E_1\) and \(E_2\) are not isogenous. By Poincar\'e's theorem,
every elliptic curve in \(A_{\overline{\Q}}\) is isogenous to \(E_1\)
or to \(E_2\), and so is the image of a nonzero homomorphism
\(E_i\to A_{\overline{\Q}}\). For \(i=1,2\), the space
\(\Hom(E_i,A_{\overline{\Q}})\otimes\Q\) has dimension one over the
division algebra \(\End^0(E_i)\). Hence all nonzero homomorphisms
\(E_i\to A_{\overline{\Q}}\) have the same image, and
\(A_{\overline{\Q}}\) contains exactly two elliptic curves.

The group \(G_{\Q}\) permutes these two elliptic curves. If it fixed
them, they would be defined over \(\Q\), and \(A\) would not be simple
over \(\Q\). Hence they are exchanged. They are defined over a quadratic
field \(K\), and the nontrivial automorphism \(\sigma\) of \(K\)
interchanges them; call them \(E\) and \(E^\sigma\).

Put \(F=A_K/E\). The composite \(E^\sigma\hookrightarrow A_K\to F\) has
finite kernel \(E\cap E^\sigma\), so it is a \(K\)-isogeny. Applying
\(\sigma\), we find that \(F^\sigma\) is \(K\)-isogenous to \(E\).

Let \(q\in\{3,7\}\). By Lemma~\ref{lem:sequence} over \(K\), either
\(P_q\in E(K)\), or the image of \(P_q\) in \(F(K)\) has order \(q\);
in the second case, its conjugate under \(\sigma\) is a point of order
\(q\) in \(F^\sigma(K)\). Thus the \(K\)-isogeny class of \(E\) contains
curves with \(K\)-rational points of orders \(3\) and \(7\). By
Lemma~\ref{lem:combine}, it contains an elliptic curve over the
quadratic field \(K\) with a point of order \(21\). This contradicts the
theorem of Kenku and Momose.

\section{The CM case: a local argument at \texorpdfstring{\(2\)}{2}}\label{sec:cm}

Suppose that \(A_{\overline{\Q}}\) is isogenous to \(E^2\), where \(E\)
has complex multiplication. The curve \(E\) is defined over a number
field and its \(j\)-invariant is an algebraic integer, so \(E\) has
potentially good reduction at every prime \cite[\S2]{ST}. Over a finite
extension of \(\Q_2\), the surface \(A\) becomes isogenous to \(E^2\)
and \(E\) acquires good reduction. By (a), \(A\) acquires good reduction
there as well, so \(A_{\Q_2}\) has potentially good reduction.

By (c), there is a finite extension \(K'/\Q_2\) with residue field
\(\F_2\) over which \(A\) has good reduction \(\widetilde A\). Since
\(63\) is odd, (b) embeds the subgroup of \(A(\Q)\subseteq A(K')\)
generated by \(P\) into \(\widetilde A(\F_2)\). Therefore
\[
 63\leq\#\widetilde A(\F_2)\leq(1+\sqrt2)^4=17+12\sqrt2<34,
\]
a contradiction. In fact,
\(\#\widetilde A(\F_2)\in\{1,\dots,16,19,20,25\}\) for every abelian surface
over \(\F_2\) (see Section~\ref{sec:computations}).
Alternatively, Clark \cite[Theorem 5]{Clark} lists the possible orders
of the rational torsion subgroup of an abelian surface over \(\Q\) with
everywhere potentially good reduction, and \(63\) is not among them.
This case does not use simplicity over \(\Q\).

\section{Non-CM squares: the endomorphism field}\label{sec:endomorphism-field}

From now on, suppose that \(A_{\overline{\Q}}\) is isogenous to \(E^2\),
where \(E\) has no complex multiplication. Then
\[
 D=\End^0(A_{\overline{\Q}})\simeq M_2(\Q).
\]
Let \(L\) be the smallest field over which all geometric endomorphisms
of \(A\) are defined. It is a finite Galois extension of \(\Q\). Put
\(\Gamma=\Gal(L/\Q)\). The group \(G_{\Q}\) acts on \(D\) by
\(g(\varphi)=g\circ\varphi\circ g^{-1}\), and this action factors
through a faithful action of \(\Gamma\). By Skolem--Noether,
\(g(\varphi)=R_g\varphi R_g^{-1}\) for some
\(R_g\in D^\times=\GL_2(\Q)\), unique up to \(\Q^\times\), and
\(g\mapsto[R_g]\) embeds \(\Gamma\) in
\(\operatorname{Aut}(M_2(\Q))=\PGL_2(\Q)\). On Tate modules,
\(g(\varphi x)=g(\varphi)\,g(x)\) for \(\varphi\in D\) and
\(x\in V_\ell(A)\).

\emph{Finite subgroups of \(\PGL_2(\Q)\).} A finite subgroup of
\(\PGL_2(\mathbb R)\) is cyclic or dihedral: it fixes a point of the
hyperbolic plane, and its action on the tangent plane embeds it in
\(O(2)\). If an element \([M]\in\PGL_2(\Q)\) has finite order \(n>2\),
then \(M\) has distinct eigenvalues whose ratio is a primitive \(n\)-th
root of unity \(\zeta\), and
\[
 \frac{(\operatorname{tr}M)^2}{\det M}=2+\zeta+\zeta^{-1}\in\Q .
\]
Hence \(n\in\{3,4,6\}\), and
\[
 \Gamma\in\{C_1,C_2,C_3,C_4,C_6,V_4,D_6,D_8,D_{12}\},
\]
where \(D_{2n}\) is the dihedral group of order \(2n\) and
\(V_4=C_2\times C_2\).

\emph{Trace-one representatives.} If \([M]\) has order
\(n\in\{3,4,6\}\), then \(\operatorname{tr}M\neq0\), and we may scale
\(M\) to a representative \(R\) with \(\operatorname{tr}R=1\). Then
\(1/\det R=2+\zeta+\zeta^{-1}=c\), where \(c=1,2,3\) for \(n=3,4,6\), so
that
\[
 R^2-R+\tfrac1c=0,\qquad
 \operatorname{disc}\bigl(X^2-X+\tfrac1c\bigr)=1-\tfrac4c=-3,\ -1,\ -\tfrac13 .
\]
Since \(R(1-R)=1/c\) is a scalar, the trace-one representative of
\([M]^{-1}\) is \(1-R\). An element \([M]\) of order \(2\) has a
representative \(R\) with \(\operatorname{tr}R=0\), and then
\(R^2=-\det R\) is a nonzero scalar.

\emph{The trivial group.} If \(\Gamma=C_1\), then \(\End^0_{\Q}(A)=D\)
contains a nontrivial idempotent, and \(A\) is not simple over \(\Q\).
Hence \(\Gamma\neq C_1\).

\begin{lemma}\label{lem:subvarieties}
Every elliptic curve \(C\subset A_{\overline{\Q}}\) is defined over
\(L\), and any two such curves are \(L\)-isogenous. If \(E\subset A_L\)
is one of them, then \(A_L/E\) is \(L\)-isogenous to \(E\).
\end{lemma}

\begin{proof}
By Poincar\'e's theorem, \(C\) has a complement, and \(C\) is the image
of \(\varphi=ne\in\End(A_{\overline{\Q}})\) for a rank-one idempotent
\(e\in D\) and an integer \(n\geq1\). Since \(\varphi\) is defined over
\(L\), so is its image. Let \(C'\) be another elliptic curve, the image
of \(\varphi'=n'e'\). Since \(D\simeq M_2(\Q)\), there is \(X\in D\)
with \(e'Xe\neq0\); choose \(k\geq1\) with
\(kX\in\End(A_{\overline{\Q}})\), and put
\(\psi=\varphi'\circ kX\circ\varphi\). Then \(\psi\) is defined over
\(L\), its image lies in \(C'\), and \(\psi\circ\varphi=n\psi\neq0\), so
\(\psi\) does not vanish on \(C\). Thus \(\psi|_C:C\to C'\) is an
\(L\)-isogeny. Finally, if \(E'\subset A_L\) is an elliptic curve with
\(E\cap E'\) finite (for instance the image of an integral multiple of
\(1-e\)), then \(E'\to A_L/E\) is an \(L\)-isogeny, and \(E'\) is
\(L\)-isogenous to \(E\).
\end{proof}

\begin{corollary}\label{cor:21}
There is an elliptic curve \(E''\) over \(L\), \(L\)-isogenous to every
elliptic curve in \(A_L\), such that \(E''(L)\) contains a point of
order \(21\).
\end{corollary}

\begin{proof}
Choose an elliptic curve \(E\subset A_L\) and put \(F=A_L/E\). By
Lemma~\ref{lem:sequence} over \(L\), for \(q\in\{3,7\}\), either
\(P_q\in E(L)\), or its image in \(F(L)\) has order \(q\). By
Lemma~\ref{lem:subvarieties}, \(F\) is \(L\)-isogenous to \(E\). So the
\(L\)-isogeny class of \(E\) contains curves with \(L\)-rational points
of orders \(3\) and \(7\), and Lemma~\ref{lem:combine} gives \(E''\). By
Lemma~\ref{lem:subvarieties}, \(E''\) is \(L\)-isogenous to every
elliptic curve in \(A_L\).
\end{proof}

\section{Three lemmas on residual representations}\label{sec:residual}

In this section \(p\) is an odd prime. Put \(V=V_p(A)\), \(T=T_p(A)\),
and \(D_p=D\otimes\Q_p\simeq M_2(\Q_p)\). The algebra \(D_p\) acts on
\(V\), and \(g(\varphi x)=g(\varphi)\,g(x)\) for \(g\in G_{\Q}\),
\(\varphi\in D_p\) and \(x\in V\). As a \(D_p\)-module, \(V\) is the sum
of two copies of the standard module \(\Q_p^2\).

\begin{lemma}[Polarization]\label{lem:polarization}
Let \(T'\subset V\) be a \(G_{\Q}\)-stable lattice such that every
composition factor of \(T'/pT'\) is trivial as a \(G_L\)-module. Then
\(\Q(\zeta_p)\subseteq L\).
\end{lemma}

\begin{proof}
A polarization of \(A\) defined over \(\Q\) and the Weil pairing give a
nondegenerate alternating \(G_{\Q}\)-equivariant pairing on \(V\) with
values in \(\Q_p(1)\), hence an isomorphism
\(V\simeq V^*\otimes\Q_p(1)\). Transport the lattice
\(\Hom_{\Z_p}(T',\Z_p)\otimes\Z_p(1)\) into \(V\) by this isomorphism.
Its residual module is \((T'/pT')^*\otimes\mu_p\), whose
\(G_L\)-composition factors are four copies of
\(\overline\chi_p|_{G_L}\). By lattice independence over \(G_L\), these
agree with the composition factors of \(T'/pT'\), which are trivial.
Hence \(\overline\chi_p|_{G_L}=1\), that is,
\(\Q(\zeta_p)\subseteq L\). The degree of the polarization plays no
role.
\end{proof}

\begin{lemma}[Unramified order]\label{lem:unramified}
Suppose \(\mathcal B\subset D_p\) is a \(\Z_p\)-order isomorphic to the
ring of integers of the unramified quadratic extension of \(\Q_p\), such
that \(R_g\mathcal BR_g^{-1}=\mathcal B\) for all \(g\in\Gamma\).
Suppose also that some \(g_0\in\Gamma\) has a representative
\(R_{g_0}\in\mathcal B^\times\) whose reduction in
\(\mathcal B/p\mathcal B\simeq\F_{p^2}\) does not lie in \(\F_p\). If
\(A(\Q)[p]\neq0\), then \(\Q(\zeta_p)\subseteq L\).
\end{lemma}

\begin{proof}
\emph{A stable maximal order.} The field \(\mathcal B\otimes\Q_p\) makes
\(\Q_p^2\) a one-dimensional vector space over it, so the
\(\mathcal B\)-stable lattices in \(\Q_p^2\) are the fractional ideals
of the discrete valuation ring \(\mathcal B\) times a fixed vector.
Since \(p\) is a uniformizer of \(\mathcal B\), they form a single
homothety class \(\{p^j\Lambda\}\). For \(g\in\Gamma\), the lattice
\(R_g\Lambda\) is stable under \(R_g\mathcal BR_g^{-1}=\mathcal B\), so
\(R_g\Lambda=p^j\Lambda\) for some \(j\). Rescaling, we choose for each
\(g\in\Gamma\) a representative \(M_g\in\GL(\Lambda)\), unique up to
\(\Z_p^\times\); for \(g_0\) we take \(M_{g_0}=R_{g_0}\), which lies in
\(\GL(\Lambda)\) because \(\mathcal B^\times\) preserves \(\Lambda\).
Thus \(\Gamma\) preserves the maximal order
\(\mathcal R=\End_{\Z_p}(\Lambda)\simeq M_2(\Z_p)\).

\emph{A tensor decomposition.} Put \(T'=\mathcal RT\). It is a lattice
in \(V\), and it is \(G_{\Q}\)-stable because
\(g(\varphi t)=g(\varphi)\,g(t)\) and \(g(\mathcal R)=\mathcal R\). By
Morita equivalence,
\[
 T'\simeq\Lambda\otimes_{\Z_p}W_0,\qquad W_0=\Hom_{\mathcal R}(\Lambda,T'),
\]
with \(W_0\) free of rank \(2\). Let \(\widetilde G\) be the group of
pairs \((g,M)\) with \(g\in G_{\Q}\), \(M\in\GL(\Lambda)\) and
\([M]=[M_{g|_L}]\) in \(\PGL(\Lambda)\). For \((g,M)\in\widetilde G\),
the map \(t\mapsto M^{-1}g(t)\) is an \(\mathcal R\)-linear automorphism
of \(T'\), hence equals \(1\otimes w\) for a unique
\(w=w(g,M)\in\GL(W_0)\). Thus \(\widetilde G\) acts linearly on
\(\Lambda\) through \(M\) and on \(W_0\) through \(w\), and \((g,M)\)
acts on \(T'\) as \(g=M\otimes w\). Reducing modulo \(p\),
\[
 \overline V:=T'/pT'\simeq U\otimes_{\F_p}W,\qquad
 U=\Lambda/p\Lambda,\quad W=W_0/pW_0,
\]
as \(\widetilde G\)-modules, and these actions factor through a finite
quotient of \(\widetilde G\). The subgroup
\(H=\{(h,1):h\in G_L\}\) is normal in \(\widetilde G\), acts trivially
on \(U\), and acts on \(\overline V\) as \(G_L\) does.

\emph{\(U\) is irreducible.} The element \((g,R_{g_0})\), for any lift
\(g\) of \(g_0\), acts on \(U\simeq\mathcal B/p\mathcal B=\F_{p^2}\) by
multiplication by an element of \(\F_{p^2}\setminus\F_p\). Its
characteristic polynomial over \(\F_p\) is irreducible, so it fixes no
line of \(U\).

\emph{\(W\) is irreducible.} Since \(A(\Q)[p]\neq0\), the module
\(\overline V\) has a trivial composition factor, as a
\(G_{\Q}\)-module and hence as a \(\widetilde G\)-module. If \(W\) had a
stable line \(\ell\), then \(\overline V\) would have the composition
factors \(U\otimes\ell\) and \(U\otimes(W/\ell)\), both irreducible of
dimension \(2\), which is impossible.

\emph{\(G_L\) acts trivially on \(\overline V\).} The restriction
\(W|_H\) is semisimple: if \(S\subseteq W\) is a simple
\(H\)-submodule, then its \(\widetilde G\)-translates are simple
\(H\)-submodules because \(H\) is normal, and their sum is a nonzero
\(\widetilde G\)-submodule, hence equal to \(W\). The module
\(\overline V|_H\simeq W|_H\oplus W|_H\) has a trivial composition
factor, so \(W^H\neq0\). Since \(H\) is normal, \(W^H\) is
\(\widetilde G\)-stable, so \(W^H=W\). Hence \(G_L\) acts trivially on
\(\overline V=U\otimes W\), and Lemma~\ref{lem:polarization} gives
\(\Q(\zeta_p)\subseteq L\).
\end{proof}

\begin{lemma}[Split blocks]\label{lem:split}
Suppose \(\mathcal B=\Z_pe_1\oplus\Z_pe_2\subset D_p\), where
\(e_1,e_2\) are orthogonal idempotents of rank one with \(e_1+e_2=1\),
is preserved by \(\Gamma\), and some element of \(\Gamma\) exchanges
\(e_1\) and \(e_2\). Let \(\Gamma_0\) be the stabilizer of \(e_1\), of
index \(2\), and \(K=L^{\Gamma_0}\). Then \(G_K\) acts on the line
\(e_1D_pe_2\) by a character \(\eta\) that factors through
\(\Gal(L/K)\); its values are roots of unity, hence lie in
\(\Z_p^\times\). If the reduction \(\overline\eta\) is nontrivial and
\(A(\Q)[p]\neq0\), then \(\Q(\zeta_p)\subseteq L\).
\end{lemma}

\begin{proof}
Put \(T'=\mathcal BT\), a \(G_{\Q}\)-stable lattice, and
\(T_i=e_iT'\), so that \(T'=T_1\oplus T_2\). The group \(G_K\) preserves
each \(T_i\), since \(g(e_it)=e_ig(t)\), while every element of
\(G_{\Q}\setminus G_K\) maps \(T_1\) onto \(T_2\) and \(T_2\) onto
\(T_1\). Put \(V_i=T_i\otimes\Q_p\) and \(M_i=T_i/pT_i\), both of
dimension \(2\).

Choose \(0\neq X\in e_1D_pe_2\). It restricts to an isomorphism
\(V_2\to V_1\), and for \(g\in G_K\) and \(v\in V_2\) we have
\(g(Xv)=\eta(g)X\,g(v)\). Hence \(V_1\simeq\eta\otimes V_2\) as
\(G_K\)-representations, and lattice independence gives
\[
 M_1^{\mathrm{ss}}\simeq\overline\eta\otimes M_2^{\mathrm{ss}}.
\]
Since \(A(\Q)[p]\neq0\), the \(G_K\)-module \(M_1\oplus M_2\) has a
trivial composition factor, so one of the \(M_i\) has one. An element
\(\tau\in G_{\Q}\setminus G_K\) identifies \(M_{3-i}\) with the
conjugate of \(M_i\) by \(\tau\), and the conjugate of a trivial
composition factor is trivial. So both \(M_1\) and \(M_2\) have a
trivial composition factor. Write
\(M_2^{\mathrm{ss}}=\mathbf 1\oplus\alpha\). Then
\(M_1^{\mathrm{ss}}=\overline\eta\oplus\overline\eta\alpha\) contains
\(\mathbf 1\), and \(\overline\eta\neq\mathbf 1\), so
\(\alpha=\overline\eta^{-1}\). All four composition factors are powers
of \(\overline\eta\), hence trivial on \(G_L\), and
Lemma~\ref{lem:polarization} gives \(\Q(\zeta_p)\subseteq L\).
\end{proof}

\section{The prime \texorpdfstring{\(7\)}{7}: the endomorphism field is cyclic of order \texorpdfstring{\(2\), \(3\) or \(6\)}{2, 3 or 6}}\label{sec:prime-7}

We apply the lemmas of Section~\ref{sec:residual} with \(p=7\), using
\(A(\Q)[7]\ni P_7\neq0\). The field \(\Q(\zeta_7)\) is cyclic of degree
\(6\) over \(\Q\). So if \(\Q(\zeta_7)\subseteq L\), then \(\Gamma\) has
a quotient isomorphic to \(C_6\). Modulo \(7\), the discriminants of the
trace-one rotations are
\[
 -3\equiv4=2^2,\qquad -1\equiv6,\qquad -\tfrac13\equiv2=3^2\pmod 7,
\]
and \(6\) is not a square modulo \(7\).

\emph{\(\Gamma=D_6\) or \(D_{12}\).} Let \(r\) be a rotation of order
\(n=3\) or \(6\) generating the rotation subgroup, and \(R\) its
trace-one representative. Every element of \(\Gamma\) conjugates \(r\)
to \(r^{\pm1}\), so it maps \(R\) to \(R\) or to \(1-R\). Hence
\(\mathcal B=\Z_7[R]\), a \(\Z_7\)-order since \(R^2-R+1/c=0\) with
\(c\in\{1,3\}\), is preserved by \(\Gamma\). Its discriminant is a
nonzero square modulo \(7\), so by Hensel's lemma the roots
\(\lambda_1,\lambda_2\) of \(X^2-X+1/c\) lie in \(\Z_7^\times\) and
\(\lambda_1-\lambda_2\) is a unit. Thus
\(\mathcal B\simeq\Z_7\times\Z_7\), with idempotents
\(e_1=(R-\lambda_2)/(\lambda_1-\lambda_2)\) and
\(e_2=(R-\lambda_1)/(\lambda_2-\lambda_1)\). The rotations fix \(R\),
hence \(e_1\) and \(e_2\). Since \(\lambda_1+\lambda_2=1\), the
reflections, which map \(R\) to \(1-R\), exchange \(e_1\) and \(e_2\).
The stabilizer of \(e_1\) is the rotation subgroup. For
\(X\in e_1D_7e_2\) we have \(RXR^{-1}=(\lambda_1/\lambda_2)X\), so
\(\eta\) maps a lift of \(r\) to \(\lambda_1/\lambda_2\), a primitive
\(n\)-th root of unity. As \(7\nmid n\), reduction modulo \(7\) is
injective on \(n\)-th roots of unity, so \(\overline\eta\) has order
\(n\geq3\). Lemma~\ref{lem:split} gives \(\Q(\zeta_7)\subseteq L\). But
the abelianizations of \(D_6\) and \(D_{12}\) are \(C_2\) and
\(C_2\times C_2\), which have no quotient \(C_6\). This is a
contradiction.

\emph{\(\Gamma=D_8\) or \(C_4\).} Let \(r\) be a rotation of order
\(4\) and \(R\) its trace-one representative, with
\(R^2-R+\tfrac12=0\). As before, every element of \(\Gamma\) maps \(R\)
to \(R\) or \(1-R\), so \(\mathcal B=\Z_7[R]\) is preserved by
\(\Gamma\). Since \(X^2-X+4\), the reduction of \(X^2-X+\tfrac12\), has
discriminant \(-15\equiv6\), a nonsquare modulo \(7\), it is irreducible
over \(\F_7\). Hence \(\mathcal B\) is the ring of integers of the
unramified quadratic extension of \(\Q_7\). Moreover
\(R\in\mathcal B^\times\), since \(\det R=\tfrac12\) is a unit, and its
reduction is a root of an irreducible quadratic, so it does not lie in
\(\F_7\). Lemma~\ref{lem:unramified} gives \(\Q(\zeta_7)\subseteq L\).
But \(C_4\) has no quotient \(C_6\), and the abelianization
\(C_2\times C_2\) of \(D_8\) has none either. This is a contradiction.

\emph{\(\Gamma=V_4\).} The map \([M]\mapsto v_7(\det M)\bmod2\) is a
well-defined homomorphism \(\Gamma\to\Z/2\Z\). Choose \(r\neq1\) in its
kernel, and a representative \(R\) of trace \(0\), scaled by a power of
\(7\) so that \(\det R\in\Z_7^\times\). Then \(R^2=a\) with
\(a=-\det R\in\Q^\times\cap\Z_7^\times\). Since \(\Gamma\) is abelian,
every \(R_g\) conjugates \(R\) to a trace-zero representative of \(r\)
with the same square, that is, to \(\pm R\). Hence
\(\mathcal B=\Z_7[R]\) is preserved by \(\Gamma\). Its discriminant
\(4a\) is a unit.
\begin{itemize}
\item If \(a\) is not a square modulo \(7\), then \(X^2-a\) is
irreducible modulo \(7\), so \(\mathcal B\) is the ring of integers of
the unramified quadratic extension of \(\Q_7\). Also
\(R\in\mathcal B^\times\) and its reduction does not lie in \(\F_7\).
Lemma~\ref{lem:unramified} applies with \(g_0=r\).
\item If \(a\) is a square modulo \(7\), then \(a=b^2\) with
\(b\in\Z_7^\times\), and \(\mathcal B\simeq\Z_7\times\Z_7\), with
idempotents \(e_\pm=(1\pm R/b)/2\). Some element of \(\Gamma\)
conjugates \(R\) to \(-R\): otherwise every \(R_g\) would lie in the
centralizer \((\Q_7[R])^\times\simeq\Q_7^\times\times\Q_7^\times\) of
\(R\), and \(\Gamma\) would embed in its image \(\Q_7^\times\) in
\(\PGL_2(\Q_7)\), whose finite subgroups are cyclic. That element
exchanges \(e_+\) and \(e_-\). The stabilizer of \(e_+\) is
\(\{1,r\}\), and \(r\) acts on \(e_+D_7e_-\) by
\(RXR^{-1}=(b/(-b))X=-X\). Thus \(\overline\eta\) is the nontrivial
quadratic character, and Lemma~\ref{lem:split} applies.
\end{itemize}
In both cases \(\Q(\zeta_7)\subseteq L\), which is impossible since
\([L:\Q]=4\).

Together with \(\Gamma\neq C_1\) (Section~\ref{sec:endomorphism-field}),
this shows:

\begin{proposition}\label{prop:cyclic}
In case \textup{(S)}, \(\Gamma\) is cyclic of order \(2\), \(3\) or
\(6\).
\end{proposition}

\section{The quadratic case}\label{sec:quadratic}

If \(\Gamma=C_2\), then \(L\) is a quadratic field, and
Corollary~\ref{cor:21} gives an elliptic curve over \(L\) with an
\(L\)-rational point of order \(21\). This contradicts the theorem of
Kenku and Momose.

\section{The cyclic sextic case: no rational \texorpdfstring{\(3\)}{3}-torsion}\label{sec:sextic}

In this section \(\Gamma=\langle r\rangle\simeq C_6\). Let \(R\) be the
trace-one representative of \(r\), so \(R^2-R+\tfrac13=0\), and let
\(K=L^{\langle r^2\rangle}\) be the quadratic subfield of \(L\). We show
that \(A[3]^{\mathrm{ss}}\) has no trivial composition factor. Since
\(P_3\) is a rational point of order \(3\), this excludes
\(\Gamma=C_6\).

\subsection{\texorpdfstring{\(L\)}{L} is totally real}\label{sec:real}

Fix an embedding \(\overline{\Q}\subset\mathbb C\), and let
\(c\in G_{\Q}\) be complex conjugation. Put
\(H=H_1(A(\mathbb C),\mathbb R)\), a real vector space of dimension
\(4\) with the complex structure \(J_A\) of the complex torus
\(A(\mathbb C)\). The algebra
\(D_{\mathbb R}=D\otimes\mathbb R\simeq M_2(\mathbb R)\) acts faithfully
on \(H\), commuting with \(J_A\). Since the simple
\(M_2(\mathbb R)\)-module is \(\mathbb R^2\), we have
\(H\simeq\mathbb R^2\otimes W\) with \(\dim W=2\), the algebra
\(D_{\mathbb R}\) acting on the first factor. The commutant of
\(M_2(\mathbb R)\otimes1\) is \(1\otimes\End(W)\), so \(J_A=1\otimes J\)
with \(J^2=-1\).

Complex conjugation on \(A(\mathbb C)\) induces an involution \(c_*\) of
\(H\) that anticommutes with \(J_A\), and the action of \(c\) on \(D\) is
\(\varphi\mapsto c_*\varphi c_*^{-1}\). By Skolem--Noether, there is
\(M\in\GL_2(\mathbb R)\) with
\(c_*(\varphi\otimes1)c_*^{-1}=M\varphi M^{-1}\otimes1\); then
\((M^{-1}\otimes1)c_*\) commutes with \(M_2(\mathbb R)\otimes1\), so
\(c_*=M\otimes N\) with \(N\in\GL(W)\). The relation
\(c_*J_A=-J_Ac_*\) gives \(NJ=-JN\). In a basis of \(W\) with
\(J=\left(\begin{smallmatrix}0&-1\\1&0\end{smallmatrix}\right)\), this
forces \(N=\left(\begin{smallmatrix}a&b\\b&-a\end{smallmatrix}\right)\),
so \(N^2=(a^2+b^2)\cdot1\) with \(a^2+b^2>0\). From \(c_*^2=1\) we get
\(M^2\otimes N^2=1\), so \(M^2=\mu\cdot1\) with
\(\mu=(a^2+b^2)^{-1}>0\). If \(M\) is not a scalar, its minimal
polynomial is \(X^2-\mu\), its eigenvalues are \(\pm\sqrt\mu\), and
\(\det M=-\mu<0\).

On the other hand, the action of \(c\) on \(D\) is \(\Ad(R_c)\) with
\([R_c]\in\Gamma\), so \(M\) is a real multiple of \(R_c\), and
\(\det M\) has the sign of \(\det R_c\); this sign is well defined on
\(\PGL_2(\mathbb R)\). Every element of \(\Gamma\) is \([R^k]\) with
\(\det R^k=3^{-k}>0\). Hence \(M\) is a scalar, \(c\) acts trivially on
\(D\), and \(c\in G_L\). Thus \(L\subset\mathbb R\), and since \(L/\Q\)
is Galois, \(L\) is totally real. In particular \(K\) is a real
quadratic field.

\subsection{An isogenous surface with an action of \texorpdfstring{\(\Z[\zeta_6]\)}{Z[zeta6]}}\label{sec:normalization}

We have \(\End^0_{\Q}(A)=D^\Gamma=\Q[R]\simeq\Q(\sqrt{-3})\), the
centralizer of \(R\). Put \(u=3R-1\). Then
\[
 u^2-u+1=0,\qquad R=\frac{1+u}3,\qquad R^2=\frac u3,
\]
so \(u\) is a primitive sixth root of unity. Put \(\mathcal O=\Z[u]\).
Choose \(m\geq1\) with \(mu\in\End_{\Q}(A)\), let
\(\iota:A\to A\times A\) be \(a\mapsto(ma,(mu)(a))\), and put
\(A'=\iota(A)\). It is \(\Q\)-isogenous to \(A\), since
\(\ker\iota\subseteq A[m]\). The endomorphism \(U(x,y)=(y,y-x)\) of
\(A\times A\) is defined over \(\Q\) and satisfies \(U^2-U+1=0\). It
preserves \(A'\): given \(a\), choose \(a''\) with \(ma''=a\); since
\((mu)^2=m(mu)-m^2\) in \(\End(A)\),
\[
 U(\iota(a))=\bigl((mu)(a),\,(mu)(a)-ma\bigr)=\iota\bigl((mu)(a'')\bigr).
\]
So \(\mathcal O\subseteq\End_{\Q}(A')\), with \(u\) acting as \(U\). The
isogeny \(A\to A'\) identifies \(D\) with \(\End^0(A'_{\overline{\Q}})\),
compatibly with the Galois action and with \(u\mapsto U\); we use the
same letters on \(A'\). By lattice independence,
\(A'[3]^{\mathrm{ss}}\simeq A[3]^{\mathrm{ss}}\).

\subsection{The twist}\label{sec:twist}

For \(g\in G_{\Q}\), write \(g|_L=r^{k(g)}\) with \(k(g)\in\Z/6\Z\), and
put \(\xi(g)=u^{-2k(g)}\). Then
\(\xi:G_{\Q}\to\mu_3=\langle u^2\rangle\subset\Aut_{\Q}(A')\) is a
homomorphism. Since its values are defined over \(\Q\), it is a
cocycle, and Galois descent gives an abelian surface \(B/\Q\) and an
isomorphism \(f:A'_{\overline{\Q}}\to B_{\overline{\Q}}\) with
\({}^gf=f\circ\xi(g)\) for all \(g\in G_{\Q}\). Transporting
endomorphisms by \(f\), the Galois action on \(D\) becomes
\[
 \varphi\longmapsto\Ad\bigl(\xi(g)R^{k(g)}\bigr)\varphi
 =\Ad\bigl(S^{k(g)}\bigr)\varphi,\qquad
 S=u^{-2}R,\quad S^2=u^{-4}\cdot\frac u3=-\frac13 .
\]
Put \(T=3S=u^{-2}(1+u)\in\mathcal O\), so that \(T^2=-3\). We deduce:
\begin{itemize}
\item \(\mathcal O\) commutes with \(\xi\), so
\(f\mathcal Of^{-1}\subseteq\End_{\Q}(B)\); we write \(u\) and \(T\)
for the corresponding endomorphisms of \(B\);
\item \(B_{\overline{\Q}}\simeq A'_{\overline{\Q}}\) is isogenous to the
square of an elliptic curve without CM;
\item since \(\Ad(S^2)=1\), all geometric endomorphisms of \(B\) are
defined over \(K\), and \({}^\tau\varphi=T\varphi T^{-1}\) for
\(\tau\in G_{\Q}\setminus G_K\);
\item for \(x\in A'(\overline{\Q})\),
\(g(f(x))=f\bigl(\xi(g)\,g(x)\bigr)\).
\end{itemize}

\subsection{A cyclic \texorpdfstring{\(3\)}{3}-isogeny}\label{sec:isogeny}

Choose a rank-one idempotent \(e\in D\) and \(n\geq1\) with
\(ne\in\End(B_K)\), and let \(E_0=ne(B)\), an elliptic curve over \(K\).
Let \(\tau\in G_{\Q}\setminus G_K\). Since
\(3\,{}^\tau(ne)=-T(ne)T\), we get \(\tau(E_0)=T(E_0)\). Moreover
\(T(E_0)\neq E_0\): otherwise \(T|_{E_0}\) would be an endomorphism of
\(E_0\) with square \(-3\), and \(E_0\) would have CM. As \(T\) is
defined over \(\Q\), the restrictions \(T|_{E_0}:E_0\to\tau E_0\) and
\(T|_{\tau E_0}:\tau E_0\to E_0\) are \(K\)-isogenies conjugate under
\(\tau\), hence of the same degree \(d\), and their composite is
\([-3]\) on \(E_0\). So \(d^2=9\), and \(T|_{E_0}\) is a cyclic
\(K\)-isogeny of degree \(3\).

\subsection{Conclusion}

\begin{lemma}\label{lem:sextic}
If \(\Gamma\simeq C_6\), then \(A[3]^{\mathrm{ss}}\) has no trivial
composition factor. In particular \(A(\Q)[3]=0\).
\end{lemma}

\begin{proof}
Suppose \(A[3]^{\mathrm{ss}}\) has a trivial composition factor.
\begin{enumerate}
\item \emph{Reduction to \(A'[T]\).} Since \(T^2=-3\), multiplication by
\(T\) maps \(A'[3]\) into \(A'[T]=\ker T\), with kernel \(A'[T]\). As
\(\#A'[3]=81\) and \(\#A'[T]=\deg T=9\), the map is onto, and it is
\(G_{\Q}\)-equivariant because \(T\) is defined over \(\Q\). Hence
\(A'[3]^{\mathrm{ss}}\simeq(A'[T]^{\mathrm{ss}})^{2}\), and
\(A'[T]^{\mathrm{ss}}\) has a trivial composition factor.
\item \emph{Passage to \(B[T]\).} In \(\mathcal O\), the elements \(T\)
and \(\zeta_3-1=u^2-1\) generate the same prime ideal above \(3\). So
\(\mu_3\) acts trivially on \(A'[T]\), and by the last item of
Section~\ref{sec:twist}, \(f:A'[T]\to B[T]\) is \(G_{\Q}\)-equivariant.
Hence \(B[T]^{\mathrm{ss}}\) has a trivial composition factor.
\item \emph{Two lines in \(B[T]\).} The group \(B[T]\) is a
two-dimensional \(\F_3\)-vector space. Put
\(G=\ker(T|_{E_0})\subset B[T]\), a \(G_K\)-stable line with character
\(\alpha\). Then \(\tau(G)=\ker(T|_{\tau E_0})\), on which \(G_K\) acts
by \(\alpha^\tau(g)=\alpha(\tau^{-1}g\tau)\). Also
\(\ker(T|_{\tau E_0})=T(E_0[3])\), since \(T(E_0[3])\) is killed by
\(T^2=-3\) and both groups have order \(3\). As \(G_K\)-modules,
\(T(E_0[3])\simeq E_0[3]/G\), and the Weil pairing gives
\(E_0[3]/G\simeq\Hom(G,\mu_3)\). Hence
\[
 \alpha^\tau=\alpha^{-1}\,\overline\chi_3|_{G_K}.
\]
\item \emph{If \(\tau G=G\),} then \(\alpha^\tau=\alpha\), so
\(\overline\chi_3|_{G_K}=\alpha^2=1\), because \(\alpha\) takes values
in \(\F_3^\times=\{\pm1\}\).
\item \emph{If \(\tau G\neq G\),} then \(B[T]=G\oplus\tau G\) as
\(G_K\)-modules, with characters \(\alpha\) and \(\alpha^\tau\). The
trivial composition factor forces one of them to be trivial, and then
both are. So \(\overline\chi_3|_{G_K}=1\).
\end{enumerate}
In both cases \(\zeta_3\in K\), that is, \(K=\Q(\sqrt{-3})\). This
contradicts Section~\ref{sec:real}.
\end{proof}

\section{The cyclic cubic case: the curve 162b1}\label{sec:cubic}

By Proposition~\ref{prop:cyclic} and Sections~\ref{sec:quadratic}
and~\ref{sec:sextic}, it remains to treat
\(\Gamma=\langle r\rangle\simeq C_3\), with \(L\) a cyclic cubic field.

\subsection{The field \texorpdfstring{\(L\)}{L} and the isogeny class}\label{sec:field}

By Corollary~\ref{cor:21}, there is an elliptic curve \(E''\) over \(L\)
and a point \(P''\in E''(L)\) of order \(21\), where \(E''\) is
\(L\)-isogenous to every elliptic curve in \(A_L\). The pair
\((E'',P'')\) defines a point \(x\in Y_1(21)(L)\), and \(Y_1(21)\) is a
fine moduli space because \(21\geq4\). By Mazur's theorem, \(x\) is not
rational, so its field of definition is \(L\) and \(x\) is a geometric
point of a closed point of degree \(3\) on \(X_1(21)\).

Let \(E_{162}\) denote the curve 162b1. Over \(\Q(\zeta_9)^+\) its
torsion subgroup is \(\Z/21\Z\) (Section~\ref{sec:computations}), so it
has exactly \(12\) points of order \(21\), forming six classes
\(\{\pm Q\}\). Since \(\Aut(E_{162})=\{\pm1\}\), these give six distinct
points of \(Y_1(21)\), each defined over \(\Q(\zeta_9)^+\) and not over
\(\Q\). They form two closed points of degree \(3\), which by
\cite[Section 6.2]{DEvHMZB} are \(D_0\) and \(D_0'\), and they are all
the geometric points of \(D_0\) and \(D_0'\). Hence \(x\) is one of
them. Therefore \(L=\Q(\zeta_9)^+\), and
\((E'',P'')\simeq(E_{162},\pm Q)\) over \(L\) for some \(Q\). In
particular, every elliptic curve in \(A_L\) is \(L\)-isogenous to
\(E_{162}\).

\subsection{Untwisting}\label{sec:untwisting}

Now \(R^2-R+1=0\), so \(u:=R\) is a primitive sixth root of unity in
\(\End^0_{\Q}(A)=D^\Gamma=\Q[R]\), and we put \(\mathcal O=\Z[u]\) and
\(F=\Q(u)\simeq\Q(\sqrt{-3})\). As in Section~\ref{sec:normalization},
we replace \(A\) by a \(\Q\)-isogenous surface \(A'\) with
\(\mathcal O\subseteq\End_{\Q}(A')\). For \(g\in G_{\Q}\), write
\(g|_L=r^{k(g)}\) with \(k(g)\in\Z/3\Z\), and put \(\xi(g)=u^{2k(g)}\).
This is a homomorphism \(G_{\Q}\to\mu_3\subset\Aut_{\Q}(A')\) with
kernel \(G_L\). As in Section~\ref{sec:twist}, it gives a twist \(B/\Q\)
and \(f:A'_{\overline{\Q}}\to B_{\overline{\Q}}\) with
\({}^gf=f\circ\xi(g)\). The Galois action on \(D\) becomes
\[
 \Ad\bigl(\xi(g)R^{k(g)}\bigr)=\Ad\bigl(u^{3k(g)}\bigr)=\Ad(\pm1)=\mathrm{id}.
\]
So every geometric endomorphism of \(B\) is defined over \(\Q\), and
\(\End^0_{\Q}(B)\simeq M_2(\Q)\). Let \(e\in\End^0_{\Q}(B)\) be a
rank-one idempotent and \(E_1\) the image of an integral multiple of
\(e\), an elliptic curve over \(\Q\). By Lemma~\ref{lem:subvarieties}
applied to \(B\), whose endomorphism field is \(\Q\), the surface \(B\)
is \(\Q\)-isogenous to \(E_1^2\).

We claim that \(E_1\) is \(\Q\)-isogenous to \(E_{162}\). Since
\(\xi|_{G_L}=1\), the isomorphism \(f\) is defined over \(L\), so
\(B_L\simeq A'_L\), and the \(\Q\)-isogeny \(A'\to A\) maps
\(f^{-1}(E_1)\) onto an elliptic curve in \(A_L\). Hence \(E_1\) is
\(L\)-isogenous to an elliptic curve in \(A_L\), and by
Section~\ref{sec:field} to \(E_{162}\). The group
\(\Hom_{\overline{\Q}}(E_1,E_{162})\) is free of rank one over \(\Z\),
because \(E_1\) has no CM. The group \(G_{\Q}\) acts on it through a
character \(\varepsilon:G_{\Q}\to\Aut(\Z)=\{\pm1\}\). This character is
trivial on \(G_L\), because there is a nonzero homomorphism defined over
\(L\). So \(\varepsilon\) factors through \(\Gal(L/\Q)\simeq C_3\), and
is therefore trivial. Hence every geometric homomorphism
\(E_1\to E_{162}\) is defined over \(\Q\), which proves the claim. In
particular \(a_p(E_1)=a_p(E_{162})\) for every prime \(p\nmid6\).

\subsection{Frobenius at \texorpdfstring{\(5\)}{5}}\label{sec:frobenius}

Since \(\End^0_{\Q}(B)\simeq M_2(\Q)\) commutes with \(G_{\Q}\), the
module \(V_\ell(B)\) is \(\Q_\ell^2\otimes eV_\ell(B)\), with
\(M_2(\Q)\) acting on the first factor, and
\(eV_\ell(B)\simeq V_\ell(E_1)\) as \(G_{\Q}\)-modules. The subfield
\(F\subset M_2(\Q)\) makes \(\Q^2\) a one-dimensional \(F\)-vector
space. Hence
\[
 V_\ell(B)\simeq(F\otimes\Q_\ell)\otimes_{\Q_\ell}V_\ell(E_1),
\]
with \(F\) acting on the first factor and \(G_{\Q}\) on the second. By
the last relation of Section~\ref{sec:twist}, transporting through \(f\),
an element \(g\in G_{\Q}\) acts on \(V_\ell(A')\) as
\(\xi(g)^{-1}\otimes\rho_{E_1}(g)\). Also \(V_\ell(A)\simeq V_\ell(A')\).

The field \(L=\Q(\zeta_9)^+\) is unramified at \(5\), so \(\xi\) is
unramified at \(5\), and so is \(V_\ell(E_1)\) for \(\ell\neq5\), since
\(E_1\) is \(\Q\)-isogenous to \(E_{162}\), which has conductor \(162\).
By (a), \(A\) has good reduction at \(5\). The prime \(5\) is inert in
\(L\) (equivalently, \(5\not\equiv\pm1\pmod9\)), so
\(\mathrm{Frob}_5\) generates \(\Gal(L/\Q)\), and
\(\zeta=\xi(\mathrm{Frob}_5)^{-1}\) is a primitive cube root of unity in
\(F\). Over \(\overline{\Q}_\ell\), the eigenvalues of
\(\mathrm{Frob}_5\) on \(V_\ell(A)\) are
\(\omega\beta_1,\omega\beta_2,\omega^{-1}\beta_1,\omega^{-1}\beta_2\),
where \(\omega\) is the image of \(\zeta\) under an embedding
\(F\to\overline{\Q}_\ell\) and \(\beta_1,\beta_2\) are the roots of
\(X^2-a_5X+5\). With \(a_5=a_5(E_{162})=0\), the characteristic
polynomial of \(\mathrm{Frob}_5\) is
\[
 P_5(X)=N_{F/\Q}\bigl(X^2-a_5\zeta X+5\zeta^2\bigr)
 =(X^2+5\zeta^2)(X^2+5\zeta)=X^4-5X^2+25 .
\]
By (b), the reduction \(\widetilde A\) of \(A\) at \(5\) satisfies
\[
 \#\widetilde A(\F_5)=P_5(1)=21 .
\]
Since \(63\) is prime to \(5\), (b) embeds
\(\langle P\rangle\simeq\Z/63\Z\) into \(\widetilde A(\F_5)\). But
\(63\) does not divide \(21\). This contradiction completes the proof of
Theorem~\ref{thm:main}.
\hfill\(\square\)

The value \(21\) depends only on the \(\Q\)-isogeny class, so no choice
of lattice, polarization or isogenous model avoids it. The bound is
attained: the census curve 21.b has torsion \(\Z/21\Z\), and its
Frobenius polynomials agree with the formula above for all primes
\(5\leq p<300\) (Section~\ref{sec:computations}).

\section{Computations and a consistency check}\label{sec:computations}

The script
\href{https://genus-2-torsion.github.io/data/knowledge/sources/qsimple_63_arithmetic.m}{\texttt{qsimple\_63\_arithmetic.m}}
checks, with an assertion for each claim, every numerical fact used
above:
\begin{enumerate}
\item the residues modulo \(7\) of Section~\ref{sec:prime-7}:
\(-3\equiv4\) and \(-\tfrac13\equiv2\) are squares, \(-1\equiv6\) is
not, and \(X^2-X+4\) and \(X^2-a\), for \(a\) a nonsquare, are
irreducible over \(\F_7\);
\item that \(\#B(\F_2)\in\{1,\dots,16,19,20,25\}\) for every abelian
surface \(B\) over \(\F_2\), by exhaustive enumeration of the
\(2\)-Weil polynomials \(X^4+aX^3+bX^2+2aX+4\) (Section~\ref{sec:cm});
\item the identities in \(\Q(\sqrt{-3})\) of Sections~\ref{sec:sextic}
and~\ref{sec:cubic}, including \(S^2=-\tfrac13\), \(T^2=-3\), and the
equality of the ideals \((T)\) and \((\zeta_3-1)\);
\item that \(\Q[x]/(x^3-3x+1)\) is the cyclic cubic field
\(\Q(\zeta_9)^+\), of discriminant \(81\), in which \(5\) is inert;
\item that 162b1 has the model above, conductor \(162\), non-integral
\(j\)-invariant \(-140625/8\) (hence no CM), and \(a_5=0\), and that its
torsion subgroup over \(\Q(\zeta_9)^+\) is \(\Z/21\Z\), with \(12\)
points of order \(21\), for instance \((1-2\alpha,2-2\alpha)\) where
\(\alpha^3-3\alpha+1=0\);
\item the polynomial \(P_5(X)=X^4-5X^2+25\) and the value \(P_5(1)=21\)
of Section~\ref{sec:cubic}.
\end{enumerate}

The script
\href{https://genus-2-torsion.github.io/data/knowledge/sources/qsimple_63_census_21b.m}{\texttt{qsimple\_63\_census\_21b.m}}
is a consistency check and is not used in the proof. The census curve
21.b,
\[
 y^2+(x^3+x+1)y=x^5+2x^4+2x^3+x^2
\]
(LMFDB 324.a.648.1) is listed in the class ``geometrically split, simple
over \(\Q\)'' with torsion \(\Z/21\Z\). The script checks that, for
every prime \(5\leq p<300\), the characteristic polynomial of Frobenius
of its Jacobian equals \((X^2-a_pX+p)^2\) if \(p\equiv\pm1\pmod9\) and
\(N_{F/\Q}(X^2-a_p\zeta X+p\zeta^2)\) otherwise, with
\(a_p=a_p(E_{162})\), as predicted by Section~\ref{sec:cubic}. It also
checks that \(\#J(\F_5)=21\) and
\(J(\Q)_{\mathrm{tors}}\simeq\Z/21\Z\). So the argument stops exactly
at \([21]\).

Both scripts were run in Magma V2.30-1, and both logs end with
\texttt{ALL\_CHECKS\_PASSED}. No step depends on the generalized Riemann
hypothesis.

\section{Consequence for the census}

The group
\href{https://genus-2-torsion.github.io/group.html?g=63}{\([63]\)} is
impossible in the class ``geometrically split, simple over \(\Q\)''. The
theorem places no restriction on the geometric splitting degree, the
endomorphism field, the integral geometric endomorphism ring, or the
polarization.

Simplicity over \(\Q\) is essential: \([63]\) has an example in the class
``split over \(\Q\)'' (the curve 63.a). Thus this exclusion does not
imply \(\emptyset\) on the combined geometrically split side, and the
combined geometrically split ``\(\infty\)?'' grade of \([63]\) is
unchanged.

\begin{thebibliography}{9}
\bibitem{KM}
M. A. Kenku and F. Momose,
\emph{Torsion points on elliptic curves defined over quadratic fields},
Nagoya Math. J. \textbf{109} (1988), 125--149, Theorem (0.1).
\href{https://doi.org/10.1017/S0027763000002816}{doi:10.1017/S0027763000002816}.

\bibitem{Mazur}
B. Mazur,
\emph{Modular curves and the Eisenstein ideal},
Publ. Math. IH\'ES \textbf{47} (1977), 33--186, Theorem (7\('\)).
\href{https://www.numdam.org/item/PMIHES_1977__47__33_0/}{Numdam}.

\bibitem{DEvHMZB}
M. Derickx, A. Etropolski, M. van Hoeij, J. S. Morrow and
D. Zureick-Brown,
\emph{Sporadic cubic torsion},
Algebra \& Number Theory \textbf{15} (2021), 1837--1864, Theorem A and
Section 6.2 (numbering of arXiv version 2).
\href{https://doi.org/10.2140/ant.2021.15.1837}{doi:10.2140/ant.2021.15.1837};
\href{https://arxiv.org/abs/2007.13929}{arXiv:2007.13929}.

\bibitem{ST}
J.-P. Serre and J. Tate,
\emph{Good reduction of abelian varieties},
Ann. of Math. (2) \textbf{88} (1968), 492--517, Theorem 1 and its
Corollary 2, Lemma 2, and \S2 (p.~496 and p.~498).
\href{https://doi.org/10.2307/1970722}{doi:10.2307/1970722}.

\bibitem{Clark}
P. L. Clark,
\emph{Bounds for torsion on abelian varieties with integral moduli},
arXiv:math/0407264v2 (2004), Theorem 5.
\href{https://arxiv.org/pdf/math/0407264}{arXiv:math/0407264}.
The same list is Theorem~1.4 of P. L. Clark and X. Xarles, \emph{Local bounds for torsion points on abelian varieties}, Canad. J. Math. \textbf{60} (2008), 532--555, \href{https://doi.org/10.4153/CJM-2008-026-x}{doi:10.4153/CJM-2008-026-x}, for abelian surfaces with everywhere anisotropic reduction, which includes potentially good reduction.
\end{thebibliography}

\end{document}
